% Pursuing Stacks (cotes 134-2 to 134-8): archive cote 134-2, site page 3 (PDF page 4).
% provisional contents. Draft compared: none.
% AI-transcribed by Claude Opus 5.5 (claude-opus-5-5, adaptive thinking, effort high), Message Batches API, batch ps1 (msgbatch_01Rv5hCaNTtUt4JHx1yqoHPL), 2026-11-12; reviewed against both supplied scan halves; see CHECK_DECISIONS.tsv.
\scanpage{134-2/3}\enlargethispage{-3pt}
\mspage{i}

\margin{\uline{NB} Les chiffres souli-gnés \uline{5}, \uline{7}, \uline{13}, \uline{App.3}, \uline{App.6}, \uline{App 11}, \uline{App 12} etc sont à imprimer en caractères \unsure{gras}}\par

TABLE OF CONTENTS of The Modilizing Story (1) .\par

I \uline{Take-off} ( A letter to ... )\par

1. The importance of innocence.\par
2. A short look into purgatory.\par
3. ``Fundamental $\infty$-groupoids'' as objects of a ``model category'' ?\par
4. A bit of ordering in the mess of ``higher order structures''\hand{.}\par
\uline{5}. Jumping over the abyss !\par
6. The topological model; hemispheres building up the (tentative)\par
``universal $\infty$-- (co)groupoid''.\par
7. Gluing hemispheres: the ``standard'' amalgamations.\par
8. Description of the universal \uline{primitive} structure.\par
\uline{9}. The main inductive step: just add coherence arrows ! And the\par
abridged story of an (unescapable and irrelevant) ambiguity ...\par
10. Cutting down redundancies - or: ``l'embarras du choix''\hand{.}\par
\unsure{\hand{\textasciicircum}}11. Back to the topological model. The canonical functor from\par
spaces to `` $\infty$- groupoids'' .\par
12. About replacing spaces by objects of a ``model category''.\par
\uline{13}. An urgent reflection on proper names: ``stacks'' and ``coherators''\hand{.}\par

\margin{espacer deux fois (with arrows to the space before and after the following line)}\par
\uline{Appendix} : Three letters to Larry Breen\par

(App.1) Exemples de 2-catégories de Picard.\par
(App.2) \xout{Lesxnxgroupoi} Théorème de Lefschetz (faible) en termes de \hand{(}Champs.\hand{)}\par
(\uline{App.3}) Les n-groupoides comme types d'homotopie tronqués.\par
(App.4) Relativisation sur un topos.\par
(App.5) Ingrédients principaux vers une ``Algèbre Topologique''.\par

\margin{espacer}\par
(\uline{App.6}) Champs essentiellement localement constants et types d'homoto\hand{(}pie\par
% NOTE: the right edge of the page is cut; the end of the App.6 line ("pie" typed above, with a hand bracket) is partly lost
(App.7) Intégration de champs et cohomologie.\par
(App.8) Les trois approches vers la cohomologie d'un topos.\par

\margin{espacer}\par
(App.9) \xout{Rsfle} Question about a non-abelian Dold-Kan theorem.\par
(App.10) The ``six operations'' and homology.\par
(\uline{App.11}) Complex of generalized jacobians.\par
(\uline{App.12}) Global ``geometric'' class field theory as a cohomological\par
duality formula. Serre duality and the ``Lang trick''.\par
(\uline{App.13}) Case of local ``geometric'' class field theory.\par
(App.14) Comments on fppf cohomology \deleted{and} \inserted{versus} \hand{crystallin cohomology.}\par
(App.15) The homotopical trinity:\deleted{\hand{(}}``spaces'', ``$\infty$-groupoids'', \hand{(}essential-\par%
-ly locally constant) ``coefficient categories''\hand{...}\par
(App.16) The\hand{``}six operations\hand{''}for P-constructible sheaves (for a given\par
equisingular stratification P).\par
(\uline{App.17}) Relation with ``indexed homotopy types'' and with ``dévissage''\par
of stratified spaces. The ``fine\hand{''} homotopy type\hand{''} of a (tame)\par
space or of a scheme.\par
(App.18) How far do\hand{``}essentially constant\hand{''} \uline{abelian} coefficients determi-\par%
ne a homotopy type ?\par

II \uline{Test Categories}\par

\unsure{\over{1}{1}}4. The unnoticed failure (of the foundations of topology).\par
15. Overall review on standard descriptions of homotopy types.\par
16. Stacks over topoi as the unifying concept for homotopical and\par
cohomological algebra.\par
16\hand{'}. Categories as models for homotopy types. First glimpse upon an\par
``impressive bunch'' (of modelizers).\par
17. The Artin-Mazur cohomological criterion for weak equivalence.\par
18. Corrections and comments to letter. Benabou's lonely approach.\par

\mspage{\hand{3}}
